% Printed 528; principal PDF 535; Google PDF 539.
\par\clearpage
\noindent\makebox[0pt][l]{\fontsize{13}{16}\selectfont 528}\hfill{\small\textbf{INHALT.}}\hfill\null\par\medskip
\begingroup
\fontsize{8.5}{11.5}\selectfont
\setlength{\parindent}{0pt}
\noindent\textit{Briefwechsel.}\par
\noindent\hbox to\linewidth{\hspace*{2em}Uber den Heliotrop: \textsc{Gauss} an \textsc{Olbers} {\fontsize{7.5}{9}\selectfont\textls[60]{1821}} Juli {\fontsize{7.5}{9}\selectfont\textls[60]{1}}\hfill.\hfill.\hfill.\hfill.\hfill.\hfill.\hfill.\hfill.\hfill.\hfill.\hfill.\hfill.\hfill.\hfill.\hfill\makebox[2.7em][c]{Seite}\makebox[1.8em][r]{{\fontsize{7.5}{9}\selectfont\textls[60]{467}}}}\par
\noindent\hbox to\linewidth{\hspace*{13em}\textsc{Gauss} an \textsc{Schumacher} {\fontsize{7.5}{9}\selectfont\textls[60]{1821}} Juli {\fontsize{7.5}{9}\selectfont\textls[60]{11}}\hfill.\hfill.\hfill.\hfill.\hfill.\hfill.\hfill.\hfill.\hfill.\hfill.\hfill.\hfill.\hfill\makebox[2.7em][c]{---}\makebox[1.8em][r]{{\fontsize{7.5}{9}\selectfont\textls[60]{470}}}}\par
\noindent\hbox to\linewidth{\hspace*{13em}\textsc{Gauss} an \textsc{Schumacher} {\fontsize{7.5}{9}\selectfont\textls[60]{1821}} Nov. {\fontsize{7.5}{9}\selectfont\textls[60]{8}}\hfill.\hfill.\hfill.\hfill.\hfill.\hfill.\hfill.\hfill.\hfill.\hfill.\hfill.\hfill.\hfill\makebox[2.7em][c]{---}\makebox[1.8em][r]{{\fontsize{7.5}{9}\selectfont\textls[60]{470}}}}\par
\noindent\hbox to\linewidth{\hspace*{13em}\textsc{Gauss} an \textsc{Schumacher} {\fontsize{7.5}{9}\selectfont\textls[60]{1822}} Juli {\fontsize{7.5}{9}\selectfont\textls[60]{10}}\hfill.\rlap{\hspace{.25em}.}\hfill.\hfill.\hfill.\hfill.\hfill.\hfill.\hfill.\hfill.\hfill.\hfill.\hfill.\hfill\makebox[2.7em][c]{---}\makebox[1.8em][r]{{\fontsize{7.5}{9}\selectfont\textls[60]{470}}}}\par
\noindent\hbox to\linewidth{\hspace*{13em}\textsc{Gauss} an \textsc{Schumacher} {\fontsize{7.5}{9}\selectfont\textls[60]{1827}} Jan. {\fontsize{7.5}{9}\selectfont\textls[60]{15}}\hfill.\hfill.\hfill.\hfill.\hfill.\hfill.\hfill.\hfill.\hfill.\hfill.\hfill.\hfill.\hfill\makebox[2.7em][c]{---}\makebox[1.8em][r]{{\fontsize{7.5}{9}\selectfont\textls[60]{471}}}}\par
\noindent\textit{Veröffentlichung.}\par
\noindent\hbox to\linewidth{\hspace*{2em}Die Berichtigung des Heliotrops\hfill.\hfill.\hfill.\hfill.\hfill.\hfill.\hfill.\hfill.\hfill.\hfill.\hfill.\hfill.\hfill.\hfill.\hfill.\hfill.\hfill.\hfill.\hfill.\hfill.\hfill.\hfill\makebox[2.7em][c]{---}\makebox[1.8em][r]{{\fontsize{7.5}{9}\selectfont\textls[60]{472}}}}\par
\noindent\textit{Nachlass.}\par
\noindent\hbox to\linewidth{\hspace*{2em}Einfluss unvollkommener Berichtigung am ersten Heliotrop\hfill.\hfill.\hfill.\hfill.\hfill.\hfill.\hfill.\hfill.\hfill.\hfill.\hfill.\hfill.\hfill\makebox[2.7em][c]{---}\makebox[1.8em][r]{{\fontsize{7.5}{9}\selectfont\textls[60]{478}}}}\par
\noindent\hbox to\linewidth{\hspace*{21.5em}für den zweiten Heliotrop\hfill.\hfill.\hfill.\hfill.\hfill.\hfill.\hfill.\hfill.\hfill.\hfill.\hfill\makebox[2.7em][c]{---}\makebox[1.8em][r]{{\fontsize{7.5}{9}\selectfont\textls[60]{479}}}}\par
\noindent\hbox to\linewidth{\hspace*{2em}Zur Berichtigung des Heliotrops\hfill.\hfill.\hfill.\hfill.\hfill.\hfill.\hfill.\hfill.\hfill.\hfill.\hfill.\hfill.\hfill.\hfill.\hfill.\hfill.\hfill.\hfill.\hfill.\hfill.\hfill.\hfill\makebox[2.7em][c]{---}\makebox[1.8em][r]{{\fontsize{7.5}{9}\selectfont\textls[60]{479}}}}\par
\noindent\hbox to\linewidth{\hspace*{2em}Bemerkungen\hfill.\hfill.\hfill.\hfill.\hfill.\hfill.\hfill.\hfill.\hfill.\hfill.\hfill.\hfill.\hfill.\hfill.\hfill.\hfill.\hfill.\hfill.\hfill.\hfill.\hfill.\hfill.\hfill.\hfill.\hfill.\hfill.\hfill.\hfill.\hfill\makebox[2.7em][c]{---}\makebox[1.8em][r]{{\fontsize{7.5}{9}\selectfont\textls[60]{483}}}}\par
\par\bigskip\begin{center}MESSUNGSFEHLER.\end{center}
\noindent\textit{Nachlass und Briefwechsel.}\par
\noindent\hbox to\linewidth{\hspace*{2em}Über die bei der Landestriangulierung erforderlichen Instrumente\hfill.\hfill.\hfill.\hfill.\hfill.\hfill.\hfill.\hfill.\hfill.\hfill.\hfill\makebox[2.7em][c]{---}\makebox[1.8em][r]{{\fontsize{7.5}{9}\selectfont\textls[60]{487}}}}\par
\noindent\hbox to\linewidth{\hspace*{2em}Uber Messungsfehler\enspace \textsc{Gauss} an \textsc{Olbers} {\fontsize{7.5}{9}\selectfont\textls[60]{1825}} Juli\hfill.\hfill.\hfill.\hfill.\hfill.\hfill.\hfill.\hfill.\hfill.\hfill.\hfill.\hfill.\hfill.\hfill.\hfill.\hfill\makebox[2.7em][c]{---}\makebox[1.8em][r]{{\fontsize{7.5}{9}\selectfont\textls[60]{490}}}}\par
\noindent\hbox to\linewidth{\hspace*{13.5em}\textsc{Gauss} an \textsc{Schumacher} {\fontsize{7.5}{9}\selectfont\textls[60]{1825}} August {\fontsize{7.5}{9}\selectfont\textls[60]{14}}\hfill.\hfill.\hfill.\hfill.\hfill.\hfill.\hfill.\hfill.\hfill.\hfill.\hfill\makebox[2.7em][c]{---}\makebox[1.8em][r]{{\fontsize{7.5}{9}\selectfont\textls[60]{493}}}}\par
\noindent\hbox to\linewidth{\hspace*{13.5em}\textsc{Gauss} an \textsc{Bessel} {\fontsize{7.5}{9}\selectfont\textls[60]{1843}} Okt. {\fontsize{7.5}{9}\selectfont\textls[60]{29}}\hfill.\hfill.\hfill.\hfill.\hfill.\hfill.\hfill.\hfill.\hfill.\hfill.\hfill.\hfill.\hfill.\hfill\makebox[2.7em][c]{---}\makebox[1.8em][r]{{\fontsize{7.5}{9}\selectfont\textls[60]{494}}}}\par
\noindent\hbox to\linewidth{\hspace*{13.5em}\textsc{Gauss} an \textsc{Bessel} {\fontsize{7.5}{9}\selectfont\textls[60]{1844}} August {\fontsize{7.5}{9}\selectfont\textls[60]{15}}\hfill.\hfill.\hfill.\hfill.\hfill.\hfill.\hfill.\hfill.\hfill.\hfill.\hfill.\hfill.\hfill\makebox[2.7em][c]{---}\makebox[1.8em][r]{{\fontsize{7.5}{9}\selectfont\textls[60]{498}}}}\par
\noindent\hbox to\linewidth{\hspace*{2em}Bemerkungen\hfill.\hfill.\hfill.\hfill.\hfill.\hfill.\hfill.\hfill.\hfill.\hfill.\hfill.\hfill.\hfill.\hfill.\hfill.\hfill.\hfill.\hfill.\hfill.\hfill.\hfill.\hfill.\hfill.\hfill.\hfill.\hfill.\hfill.\hfill.\hfill\makebox[2.7em][c]{---}\makebox[1.8em][r]{{\fontsize{7.5}{9}\selectfont\textls[60]{499}}}}\par
\noindent\textit{Briefwechsel.}\par
\noindent\hspace*{2em}Berechnung des mittleren Ablesungsfehlers einschließlich des Teilungsfehlers am Theodo-\par
\noindent\hbox to\linewidth{\hspace*{4.5em}liten: \textsc{Gauss} an \textsc{Gerling} {\fontsize{7.5}{9}\selectfont\textls[60]{1844}} April {\fontsize{7.5}{9}\selectfont\textls[60]{17}}\hfill.\hfill.\hfill.\hfill.\hfill.\hfill.\hfill.\hfill.\hfill.\hfill.\hfill.\hfill.\hfill.\hfill.\hfill.\hfill.\hfill\makebox[2.7em][c]{---}\makebox[1.8em][r]{{\fontsize{7.5}{9}\selectfont\textls[60]{500}}}}\par
\noindent\textit{Nachlass.}\par
\noindent\hbox to\linewidth{\hspace*{2em}Berechnung des mittleren Ablesungsfehlers am Theodoliten und \textsc{Borda}schen Kreis {\fontsize{7.5}{9}\selectfont\textls[60]{1821---1824}}\hfill\makebox[2.7em][c]{---}\makebox[1.8em][r]{{\fontsize{7.5}{9}\selectfont\textls[60]{501}}}}\par
\noindent\hbox to\linewidth{\hspace*{2em}Bemerkungen\hfill.\hfill.\hfill.\hfill.\hfill.\hfill.\hfill.\hfill.\hfill.\hfill.\hfill.\hfill.\hfill.\hfill.\hfill.\hfill.\hfill.\hfill.\hfill.\hfill.\hfill.\hfill.\hfill.\hfill.\hfill.\hfill.\hfill.\hfill.\hfill\makebox[2.7em][c]{---}\makebox[1.8em][r]{{\fontsize{7.5}{9}\selectfont\textls[60]{505}}}}\par
\par\bigskip\begin{center}REDUKTION SCHIEFER WINKEL AUF DEN HORIZONT.\quad REPETITIONSBEOBACHTUNGEN.\end{center}
\noindent\textit{Nachlass.}\par
\noindent\hbox to\linewidth{\hspace*{2em}Reduktion schiefer Winkel auf den Horizont\hfill.\hfill.\hfill.\hfill.\hfill.\hfill.\hfill.\hfill.\hfill.\hfill.\hfill.\hfill.\hfill.\hfill.\hfill.\hfill.\hfill.\hfill\makebox[2.7em][c]{---}\makebox[1.8em][r]{{\fontsize{7.5}{9}\selectfont\textls[60]{509}}}}\par
\noindent\hbox to\linewidth{\hspace*{2em}Bemerkungen\hfill.\hfill.\hfill.\hfill.\hfill.\hfill.\hfill.\hfill.\hfill.\hfill.\hfill.\hfill.\hfill.\hfill.\hfill.\hfill.\hfill.\hfill.\hfill.\hfill.\hfill.\hfill.\hfill.\hfill.\hfill.\hfill.\hfill.\hfill.\hfill\makebox[2.7em][c]{---}\makebox[1.8em][r]{{\fontsize{7.5}{9}\selectfont\textls[60]{512}}}}\par
\noindent\textit{Nachlass und Briefwechsel.}\par
\noindent\hbox to\linewidth{\hspace*{2em}Repetitionsbeobachtungen\hfill.\hfill.\hfill.\hfill.\hfill.\hfill.\hfill.\hfill.\hfill.\hfill.\hfill.\hfill.\hfill.\hfill.\hfill.\hfill.\hfill.\hfill.\hfill.\hfill.\hfill.\hfill.\hfill.\hfill.\hfill\makebox[2.7em][c]{---}\makebox[1.8em][r]{{\fontsize{7.5}{9}\selectfont\textls[60]{514}}}}\par
\noindent\hbox to\linewidth{\hspace*{2em}Über die Fehler bei Repetitionsbeobachtungen\enspace \textsc{Gauss} an \textsc{Bessel} {\fontsize{7.5}{9}\selectfont\textls[60]{1819}} Jan. {\fontsize{7.5}{9}\selectfont\textls[60]{27}}\hfill.\hfill.\hfill.\hfill.\hfill\makebox[2.7em][c]{---}\makebox[1.8em][r]{{\fontsize{7.5}{9}\selectfont\textls[60]{515}}}}\par
\noindent\hbox to\linewidth{\hspace*{2em}Bemerkungen\hfill.\hfill.\hfill.\hfill.\hfill.\hfill.\hfill.\hfill.\hfill.\hfill.\hfill.\hfill.\hfill.\hfill.\hfill.\hfill.\hfill.\hfill.\hfill.\hfill.\hfill.\hfill.\hfill.\hfill.\hfill.\hfill.\hfill.\hfill.\hfill\makebox[2.7em][c]{---}\makebox[1.8em][r]{{\fontsize{7.5}{9}\selectfont\textls[60]{516}}}}\par
\begin{center}\rule{2cm}{.3pt}\end{center}
\noindent\hbox to\linewidth{\hspace*{0pt}\textit{Bemerkungen zum neunten Band}\hfill.\hfill.\hfill.\hfill.\hfill.\hfill.\hfill.\hfill.\hfill.\hfill.\hfill.\hfill.\hfill.\hfill.\hfill.\hfill.\hfill.\hfill.\hfill.\hfill.\hfill.\hfill.\hfill\makebox[2.7em][c]{---}\makebox[1.8em][r]{{\fontsize{7.5}{9}\selectfont\textls[60]{521}}}}\par
\vspace{1.3cm}
\begin{center}
\rule{4cm}{.4pt}\par\bigskip
{\fontsize{7.5}{9}\selectfont Göttingen, Druck der Dieterichschen Univ.-Buchdruckerei (W. Fr. Kaestner).}
\end{center}
\endgroup

% Principal PDF 536: unnumbered blank leaf. No asserted equivalence to a Google leaf.
\par\clearpage
\thispagestyle{empty}\null

% Principal PDF 537: unnumbered publisher advertisement; no counterpart in Google 540–544.
\par\clearpage
\thispagestyle{empty}
\begin{tikzpicture}[remember picture,overlay]
\draw[line width=.7pt] ([xshift=8mm,yshift=-9mm]current page.north west) rectangle ([xshift=-8mm,yshift=9mm]current page.south east);
\draw[line width=.35pt] ([xshift=16mm,yshift=-17mm]current page.north west) rectangle ([xshift=-16mm,yshift=17mm]current page.south east);
\end{tikzpicture}
\begingroup
\fontsize{9}{12}\selectfont
\vspace*{1.3cm}
\begin{center}
{\fontsize{18}{22}\selectfont\textls[90]{C. F. GAUSS WERKE}}\par\bigskip
{\fontsize{11}{14}\selectfont HERAUSGEGEBEN VON DER KGL. GESELLSCHAFT DER WISSENSCHAFTEN\par
ZU GÖTTINGEN.}\par\medskip
{\fontsize{10}{13}\selectfont IN KOMMISSION BEI B. G. TEUBNER IN LEIPZIG.}\par\medskip
\rule{3cm}{.3pt}\par\bigskip
{\fontsize{13}{16}\selectfont A.\quad AUSGABE AUF DRUCKPAPIER.}
\end{center}
\noindent
\begin{tabular}{@{}l@{\hspace{1em}}r@{\hspace{.8em}}p{.83\linewidth}@{}}
Band & I. & DISQUISITIONES ARITHMETICAE. Zweiter Abdruck. {\fontsize{8}{10}\selectfont 1870.} Preis Mk. {\fontsize{8}{10}\selectfont 20.---.}\\
& II. & HÖHERE ARITHMETIK. Zweiter Abdruck. {\fontsize{8}{10}\selectfont 1876.} Preis Mk. {\fontsize{8}{10}\selectfont 20.---.}\\
& & Nachträge zum ersten Abdruck des zweiten Bandes. {\fontsize{8}{10}\selectfont 1876.} Preis Mk. {\fontsize{8}{10}\selectfont 2.---.}\\
& III. & ANALYSIS. Zweiter Abdruck. {\fontsize{8}{10}\selectfont 1876.} Preis Mk. {\fontsize{8}{10}\selectfont 20.---.}\\
& IV. & WAHRSCHEINLICHKEITSRECHNUNG UND GEOMETRIE. Zweiter Abdruck. {\fontsize{8}{10}\selectfont 1880.}\newline Preis Mk. {\fontsize{8}{10}\selectfont 25.---.}\\
& V. & MATHEMATISCHE PHYSIK. Zweiter Abdruck. {\fontsize{8}{10}\selectfont 1877.} Preis Mk. {\fontsize{8}{10}\selectfont 25.---.}\\
& VI. & ASTRONOMISCHE ABHANDLUNGEN UND AUFSÄTZE. {\fontsize{8}{10}\selectfont 1874.} Preis Mk. {\fontsize{8}{10}\selectfont 33.---.}\\
& VIII. & ARITHMETIK, ANALYSIS, WAHRSCHEINLICHKEITSRECHNUNG, GEOMETRIE.\newline (Nachträge zu Band I---IV.) {\fontsize{8}{10}\selectfont 1900.} Preis Mk. {\fontsize{8}{10}\selectfont 24.---.}\\
& IX. & GEODÄSIE. (Fortsetzung von Band IV.) {\fontsize{8}{10}\selectfont 1903.} Preis Mk. {\fontsize{8}{10}\selectfont 26.---.}
\end{tabular}
\par\medskip\noindent Die Bände VII und X nebst einem Generalregister befinden sich in Vorbereitung.
\begin{center}
{\fontsize{13}{16}\selectfont\mbox{\llap{\raisebox{.4ex}{\fontsize{5}{6}\selectfont,}\hspace{2.5em}}B.\quad AUSGABE AUF SCHREIBPAPIER.}}
\end{center}
\noindent Band I---VI (nur vollständig abzugeben). Preis Mk. {\fontsize{8}{10}\selectfont 180.---.}\par
\noindent Band VIII. Preis Mk. {\fontsize{8}{10}\selectfont 30.---.}\par
\noindent Band IX. Preis Mk. {\fontsize{8}{10}\selectfont 32.---.}\par
\begin{center}\rule{3cm}{.3pt}\end{center}
\par\noindent\hspace*{\parindent}{\fontsize{11}{14}\selectfont Berichte über den Stand der Herausgabe von \textsc{Gauss}' Werken werden, von 1898 beginnend, alljährlich in den Nachrichten der Kgl. Gesellschaft der Wissenschaften zu Göttingen (geschäftliche Mitteilungen) veröffentlicht.\par}
\vfill
\begin{center}
\rule{8cm}{.3pt}\par\bigskip
{\fontsize{14}{18}\selectfont\textls[400]{GÖTTINGEN}}\par\bigskip
{\fontsize{8}{11}\selectfont\textbf{GEDRUCKT IN DER DIETERICHSCHEN UNIVERSITÄTS-BUCHDRUCKEREI}}\par\medskip
{\fontsize{8}{11}\selectfont\textbf{W. FR. KAESTNER.}}\par\bigskip
\rule{8cm}{.3pt}
\end{center}
\endgroup
